\documentclass[11pt,a4paper]{article}

% ---------------------------------------------------------
% PREAMBLE (stable, warning-free, Overleaf-safe)
% ---------------------------------------------------------
\usepackage[utf8]{inputenc}
\usepackage[T1]{fontenc}
\usepackage{lmodern}
\usepackage{amsmath, amssymb, amsfonts}
\usepackage{bm}
\usepackage{geometry}
\usepackage{graphicx}
\usepackage{hyperref}
\usepackage{cite}
\usepackage{authblk}
\usepackage{physics}
\usepackage{mathtools}

\geometry{margin=2.4cm}

\title{\textbf{Companion XLVI: The CMB Lensing Potential in the Relaxation-Driven Cyclic Cosmology}}
\author{Michael Lehmann \\ \texttt{mi.lehmann@gmx.de}}
\date{April 2026}

\begin{document}
\maketitle

\begin{abstract}
CMB lensing probes the integrated gravitational potential between the surface of 
last scattering and the observer. It encodes the geometry of spacetime, the growth 
of structure, and the distribution of matter across cosmic time. In the standard 
cosmological model, the lensing potential is determined by the sum of the two 
Newtonian potentials, which evolve identically in the absence of anisotropic 
stress. In the Relaxation-Driven Cyclic Cosmology (RDCC), the scale-dependent 
evolution of the gravitational potential modifies the deflection field, the 
mode-coupling structure of the CMB, and the lensing power spectrum. This Companion 
develops a continuous description of CMB lensing in RDCC, deriving how the 
relaxation scale influences the lensing kernel, the deflection angle, and the 
reconstruction of the lensing potential. The resulting framework provides a 
distinctive observational signature of the RDCC gravitational field.
\end{abstract}
% ---------------------------------------------------------
\section{Introduction}

CMB lensing arises from the deflection of cosmic microwave background photons by 
the gravitational potential of large-scale structure. The resulting distortions 
encode the integrated gravitational field along the line of sight and provide a 
powerful probe of the geometry and growth of structure.

In the standard cosmological model, the lensing potential is determined by the sum 
of the two Newtonian potentials, which evolve identically in the absence of 
anisotropic stress. In RDCC, the gravitational potential evolves differently. The 
relaxation mechanism suppresses the decay of small-scale modes and enhances the 
coherence of large-scale modes. This modifies the lensing potential, the deflection 
field, and the mode-coupling structure of the CMB in a scale-dependent manner. CMB 
lensing becomes a direct probe of the relaxation scale $k_{\mathrm{rel}}$ and the 
temporal structure of the RDCC gravitational field.
% ---------------------------------------------------------
\section{The RDCC Lensing Potential}

The CMB lensing potential is defined as
\begin{equation}
    \phi(\hat{\mathbf{n}})
    = -2 \int_{0}^{\chi_{*}}
      d\chi\,
      \frac{\chi_{*} - \chi}{\chi_{*}\chi}\,
      \left[\Phi + \Psi\right](\chi,\hat{\mathbf{n}}),
\end{equation}
where $\chi_{*}$ is the comoving distance to the surface of last scattering.

In RDCC, the metric potentials evolve differently:
\begin{align}
    \Phi_{\mathrm{RDCC}}(k,a)
    &= \Phi(k,a)
       \left[
         1 - \chi_{\Phi}
         \left(\frac{k}{k_{\mathrm{rel}}}\right)^{\alpha}
       \right], \\
    \Psi_{\mathrm{RDCC}}(k,a)
    &= \Psi(k,a)
       \left[
         1 - \chi_{\Psi}
         \left(\frac{k}{k_{\mathrm{rel}}}\right)^{\alpha}
       \right].
\end{align}

The RDCC lensing potential becomes
\begin{equation}
    \phi_{\mathrm{RDCC}}
    = \phi
      \left[
        1 - \chi_{\phi}
        \left(\frac{k}{k_{\mathrm{rel}}}\right)^{\alpha}
      \right].
\end{equation}

This modification suppresses small-scale lensing power and enhances large-scale 
coherence.
% ---------------------------------------------------------
\section{Deflection Fields and Mode Coupling}

The deflection angle is given by
\begin{equation}
    \mathbf{d}(\hat{\mathbf{n}})
    = \nabla \phi(\hat{\mathbf{n}}).
\end{equation}

In RDCC, the modified lensing potential produces a deflection field with smoother 
small-scale structure and enhanced large-scale coherence. The mode-coupling 
structure of the CMB becomes
\begin{equation}
    \tilde{X}_{\ell m}
    = X_{\ell m}
      + \sum_{\ell' m'}
        \sum_{LM}
        \phi_{LM}\,
        \mathcal{G}_{\ell \ell' L}^{m m' M}\,
        X_{\ell' m'},
\end{equation}
where $\mathcal{G}$ is a geometric coupling coefficient.

Because $\phi_{LM}$ is modified in RDCC, the mode-coupling structure acquires a 
distinctive scale dependence that reflects the relaxation scale.
% ---------------------------------------------------------
\section{The Lensing Power Spectrum}

The lensing potential power spectrum is defined as
\begin{equation}
    C_{L}^{\phi\phi}
    = \langle \phi_{LM} \phi^{*}_{LM} \rangle.
\end{equation}

In RDCC, this spectrum becomes
\begin{equation}
    C_{L}^{\phi\phi,\mathrm{RDCC}}
    = C_{L}^{\phi\phi}
      \left[
        1 - \chi_{\phi\phi}
        \left(\frac{L}{L_{\mathrm{rel}}}\right)^{\alpha}
      \right],
\end{equation}
where $L_{\mathrm{rel}}$ corresponds to the angular scale of the relaxation scale 
$k_{\mathrm{rel}}$.

This modification suppresses small-scale lensing power and enhances large-scale 
structure, producing a distinctive RDCC signature.
% ---------------------------------------------------------
\section{Reconstruction of the Lensing Potential}

CMB lensing reconstruction uses quadratic estimators to extract the lensing 
potential from mode coupling. The standard estimator is
\begin{equation}
    \hat{\phi}_{LM}
    = N_{L}
      \sum_{\ell m}
      \sum_{\ell' m'}
      X_{\ell m}
      X_{\ell' m'}
      \mathcal{W}_{\ell \ell' L}^{m m' M},
\end{equation}
where $N_{L}$ is a normalization factor.

In RDCC, the modified mode-coupling structure alters the reconstruction noise and 
the estimator response. The reconstructed potential becomes
\begin{equation}
    \hat{\phi}_{LM,\mathrm{RDCC}}
    = \hat{\phi}_{LM}
      \left[
        1 - \chi_{\mathrm{rec}}
        \left(\frac{L}{L_{\mathrm{rel}}}\right)^{\alpha}
      \right].
\end{equation}

This modification provides a direct observational probe of the relaxation scale.
% ---------------------------------------------------------
\section{Conclusion}

CMB lensing in the Relaxation-Driven Cyclic Cosmology reflects the scale-dependent 
evolution of the gravitational potential. The relaxation mechanism modifies the 
lensing potential, the deflection field, and the lensing power spectrum, producing 
distinctive signatures in the mode-coupling structure of the CMB. These effects 
provide a direct observational window into the relaxation scale and the temporal 
structure of the RDCC gravitational field. The present Companion establishes the 
theoretical foundation for future studies of CMB lensing tomography, lensing 
reconstruction, and the interplay between matter and light in RDCC.

% ========================
% References
% ========================

\begin{thebibliography}{10}

\bibitem{Flagship} Lehmann, M. (2026). Relaxation-Driven Cyclic Cosmology (RDCC): A Minimal CPT-Symmetric Framework with a Single Parameter. Flagship Paper. Zenodo: \url{https://zenodo.org/records/19744958} (v43.0 or later).

\bibitem{Zenodo} The complete RDCC companion series (I--LVII) is available at the same Zenodo record above.

% Thematisch relevante Companions
\bibitem{CompXVI} Companion XVI: Boltzmann Code and Modified Hierarchy.
\bibitem{CompXXXI} Companion XXXI: RDCC Halo Model and Non-Linear Structure.
\bibitem{CompXXXII} Companion XXXII--XXXIX: Galaxy--Halo Connection, Cosmic Web, Lensing, RSD, ISW, tSZ, Rotation Curves, CGM.

% Wichtige externe Referenzen
\bibitem{Planck2020} Planck Collaboration (2020), Astron. Astrophys. 641, A6.
\bibitem{DESI2024} DESI Collaboration (2024/2025), arXiv: [aktuelle $f\sigma_8$/BAO-Paper].
\bibitem{JWST2024} Maiolino et al. (2024), Nature (JWST high-$z$ galaxies/quasars).
\bibitem{Kaiser1987} Kaiser, N. (1987), Mon. Not. R. Astron. Soc. 227, 1 (RSD).
\bibitem{Bartelmann2001} Bartelmann, M. \& Schneider, P. (2001), Phys. Rep. 340, 291 (Weak Lensing Review).
\bibitem{HuOkamoto2002} Hu, W. \& Okamoto, T. (2002), Astrophys. J. 574, 566 (CMB Lensing).
\bibitem{LewisChallinor2006} Lewis, A. \& Challinor, A. (2006), Phys. Rep. 429, 1 (CMB Lensing Review).
\bibitem{NFW1997} Navarro, J. F., Frenk, C. S. \& White, S. D. M. (1997), Astrophys. J. 490, 493 (Halo Profiles).

\end{thebibliography}

\end{document}
